\chapter{相关理论与关键技术}

\section{图像二分类}

设输入胸片为$x$，真实标签$y\in\{0,1\}$，其中0表示阴性类，1表示肺炎类。模型$f_\theta$输出两个logit $z_0,z_1$，softmax将其转换为概率：
\begin{equation}
p(y=k\mid x)=\frac{\exp(z_k)}{\exp(z_0)+\exp(z_1)}.
\end{equation}
默认决策以$p(y=1\mid x)\ge 0.5$判为肺炎。概率排序、固定阈值分类和概率可信度是三个不同层面：AUROC关心排序，混淆矩阵指标依赖阈值，Brier score和ECE关心概率与实际结果的一致程度。

\section{卷积网络与Transformer}

DenseNet通过密集连接复用前面各层特征，改善梯度传播并控制参数规模\cite{huang2017densenet}。胸片中的局部纹理、肺野结构和边缘变化适合卷积特征提取，因此DenseNet121被选为主CNN基线。ConvNeXt在卷积框架内吸收分层设计、大卷积核和现代训练经验，用作现代CNN对照\cite{liu2022convnext}。

ViT把图像划分为固定大小patch，将patch嵌入与位置编码送入Transformer编码器，通过自注意力建模较长距离关系\cite{dosovitskiy2021vit}。ViT参数较多，对预训练和学习率敏感。本项目曾因配置文件学习率与发布权重真实设置不一致，得到明显欠拟合结果。这一事件说明，架构名称不能代替完整训练配方；模型比较必须同时披露初始化、优化器、学习率、batch size、epoch和checkpoint选择规则。

\section{迁移学习与类别不均衡}

训练集规模不足以支持所有模型稳定从零学习，因此三模型均从ImageNet权重初始化。迁移学习并不意味着自然图像特征与胸片完全匹配，只表示提供较好的参数起点。训练集肺炎图像多于阴性图像，普通交叉熵可能偏向多数类。加权交叉熵写为
\begin{equation}
\mathcal{L}_{CE}=-w_y\log p(y\mid x),
\end{equation}
其中$w_y$与类别频数反比。类别权重可改善少数类学习，也可能改变概率校准和固定阈值行为，因此必须同时报告召回率、特异度和Brier score。

标签平滑把硬标签替换为接近0和1的软目标，抑制过度自信。MixUp则对图像和标签做凸组合\cite{zhang2018mixup}。本项目正式低成本方案采用轻度亮度/对比度扰动与0.1标签平滑，目的不是构造复杂新模型，而是检验减少风格依赖能否改善外部固定阈值结果。

\section{患者级数据泄漏}

医学影像的统计单位通常是患者，而不是图像。同一患者的多次拍摄包含稳定的解剖、体型和采集信息。若以图像为单位随机划分，同一患者可能跨越训练和测试，独立同分布假设被破坏。即使文件内容不完全相同、标签也一致，模型仍可能借助患者特征降低测试难度。因此，数据审计必须同时检查精确文件重复和患者标识交叉。

患者级划分也不能自动保证外部有效性。一个数据集内部的患者虽然独立，采集医院、人群和标签规则仍可能相同。患者级内部评价回答“对同来源未见患者是否有效”，外部评价回答“面对来源变化是否仍有效”，两者不能互相替代。

\section{域偏移与数据捷径}

域偏移指训练与评价数据的联合分布不同。胸片中常见来源包括设备、视图、年龄、灰度映射、压缩方式、边框和疾病定义。Zech等发现，肺炎模型跨医院的泛化表现具有明显差异\cite{zech2018}。模型还可能利用与标签偶然相关的来源线索，这类线索在开发集上有效，却不会稳定迁移。

本文用两类证据检查域偏移。一类是直接比较内部和RSNA外部指标；另一类是只用简单图像统计训练来源分类器。如果均值、分位数、边缘亮度和中心纹理就能区分来源，则至少说明数据风格差异明显。高来源分类AUC并不能证明某个具体因素导致肺炎误判，但能够否定“两个数据集除疾病标签外基本一致”的假设。

混合训练把多个来源直接合并；域均衡采样让每个来源贡献近似相同的采样概率；CORAL等方法进一步对齐特征二阶统计\cite{sun2016coral}。越复杂的方法越需要严格验证，否则可能通过目标域开发数据获得局部收益，却损害未见来源。

\section{判别、校准与操作点}

准确率受类别比例影响。平衡准确率定义为召回率与特异度的平均：
\begin{equation}
\mathrm{BAcc}=\frac{1}{2}(\mathrm{Sensitivity}+\mathrm{Specificity}).
\end{equation}
F1强调阳性类precision和recall，不能反映阴性误报全貌。AUROC衡量排序，AUPRC更受阳性率影响。本文同时报告多项指标，避免单一数字替代完整误差结构。

Brier score为概率与二值结果的均方误差：
\begin{equation}
\mathrm{Brier}=\frac{1}{N}\sum_{i=1}^N(p_i-y_i)^2.
\end{equation}
ECE把样本按置信度分箱，计算箱内准确率与平均置信度差的加权和。ECE依赖分箱方式，小样本时尤其不稳定。温度缩放用单一正参数$T$缩放logit，并在验证集上最小化负对数似然\cite{guo2017calibration}。它通常不改变排序，但会改变概率和固定阈值结果。

\section{模型解释的边界}

Grad-CAM根据目标类别梯度对卷积特征图加权，得到输入区域对当前预测的贡献热图\cite{selvaraju2017gradcam}。热图能帮助检查模型是否关注图像边缘或胸腔区域，却不提供病灶真值。没有框或分割标注时，热区不能被称为病灶定位；即使与RSNA框重叠，也只能作粗粒度一致性分析。

\section{报告规范}

TRIPOD+AI强调完整报告数据来源、参与者、结局、模型构建、内部评价、性能指标和不确定性\cite{collins2024tripod}；PROBAST+AI用于检查预测模型研究的偏倚和适用性风险\cite{moons2025probast}。本报告是课程设计，不宣称遵循临床研究注册流程，但借鉴其透明性要求，区分开发数据、内部测试和外部评价，并报告失败结果与限制。

\section{为什么高内部准确率仍可能不可靠}

内部准确率同时受任务可分性、数据重复、标签噪声和划分方式影响。若训练与测试来自相同采集流程，模型可能学习一组在该流程内稳定、在外部流程中无效的相关特征。例如，某类图像更常出现特定裁剪范围，模型可以利用胸腔外背景间接预测标签。它在内部test上仍是“真正未见图像”，却没有学习预期医学构念。

评价可靠性至少包含三层。第一层是数据独立性：同一患者和同一文件不能跨集合。第二层是方案独立性：test不能反复参与超参数和模型选择。第三层是来源独立性：至少一个外部数据应来自不同采集与标注过程。三层缺一时，结论需要相应降级，而不是简单标记为“有泄漏/无泄漏”。

本项目的新患者级test解决第一层，在本轮开始时也相对旧方案保持新鲜；随着多个新方案被评价，第二层保护会逐渐减弱。RSNA提供第三层差异，但又参与混合训练，混合方案的RSNA test不再是未知域验证。报告逐项说明这些状态，避免一个“外部验证”标签掩盖不同问题。

\section{评价指标之间为何会冲突}

设模型A把所有阳性排在阴性之前，但概率普遍接近1；模型B排序稍差，却在0.5附近更平衡。A可以有更高AUC、更差Brier和更低specificity；B则可能在默认阈值有更高balanced accuracy。低成本稳健DenseNet正体现这种冲突：RSNA固定阈值改善，AUC下降。

指标冲突不是评价失败，而是模型行为具有多个维度。课程报告需要先说明用途，再决定主指标。本项目没有临床成本函数，采用balanced accuracy作为双类操作点门槛，同时保留AUC、Brier和ECE作为否决信息。如果balanced accuracy提高而AUC明显下降，结论限定为操作点改善。

阈值选择还会改变precision与recall，但不改变原始概率排序。温度缩放一般保持AUC，却会通过0.5边界改变分类标签。将“校准后accuracy变化”解释成模型重新学到更好特征是不准确的，它只是输出映射变化。

\section{课程设计中的负结果}

负结果不是没有完成工作。一个可复核的负结果需要明确假设、唯一主要变量、相同预算、完整预测和限制解释。旧全模型微调、低成本稳健方案AUC下降、域均衡未超过简单混合，都回答了具体问题。删除这些结果会产生选择性报告，使最终方案看起来比实际更确定。

本项目采用停止规则约束方法数量：简单混合达到双域门槛后，不继续使用RSNA test挑选更复杂方法。该规则牺牲了获得更高单次数字的机会，却提高研究解释的可信度。这也是模式识别实验设计的一部分。
